\section{Experimental Setup}
\label{sec:setup}

\paragraph{Corpus.} 93 labeled incidents (43 Train-Ticket, 50 Sock~Shop) form the
working set; a 23-incident \emph{gate} (one per fault family per application) bounds
agent-cost sweeps. Non-LLM baselines are additionally swept over the full
$93\times15$ condition grid.

\paragraph{Methods.} (1)~a statistical rule tree over the same tools; (2)~RCAEval
multi-source BARO~\cite{pham2024baro,pham2025rcaeval} via its published artifact;
(3)~MicroRank~\cite{yu2021microrank} and TraceRCA~\cite{li2021tracerca} (trace-only);
(4)~\agent. All methods receive identical incident windows; only the agent is subject
to masking (which, if anything, disadvantages it --- the statistical baseline still
exploits container names).

\paragraph{Metrics.} Top-1 service (AC@1), fault-type accuracy against a closed
13-label taxonomy, their conjunction (\emph{both}); AC@3/MRR for ranking baselines;
per-diagnosis tool calls, bytes touched, tokens (with cache split), and dollar cost.

\paragraph{Degradation conditions.} Seeded, deterministic, offline transforms applied
before the tool layer; the agent is held fixed within any sweep. Axes: trace retention
(100/50/25/10/5\% whole-trace sampling), metric resolution (5/10/30/60\,s), log level
(ALL/WARN/ERROR), kernel tier (all / L1-only / L3-only / none), whole-modality
removal. \todo{metric/log axis numbers; kernel$\times$trace interaction grid}

\paragraph{Reproducibility.} Every result row is keyed by (run, condition spec, seed,
method) and linked to its transcript; artifacts ship as sha256-manifested bundles.
\todo{model/version/hardware box}
